\begin{resumo}[Abstract]
\begin{otherlanguage*}{english}
This qualifying text examines land use and land cover dynamics in the state of
Goiás, Brazil, between 1985 and 2024, and their relationship with socioeconomic
drivers, using the annual MapBiomas series (Collection 10.1) for the 246
municipalities of Goiás and a mesh of 166 Minimum Comparable Areas that keeps
territory constant in longitudinal analyses. A data-driven periodization
divides the 40 years into three acts --- pasture as inheritance (1985--2000),
expansion and intensification (2001--2019), and accelerated conversion
(2020--2024) --- and frames the central finding: a northward spatial
reorganization of land conversion within the state (``the northward march'').
The investigation unfolds along four fronts: the existence of the pattern
(displacement of the conversion center of mass, controlling for scale
effects); the local mechanism (two simultaneous logics --- intensification and
pasture-to-cropland substitution in the South, expansion over natural
vegetation in the North); a formal test of intra-state causal displacement,
which is not confirmed and points to an alternative reading as a coordinated
response to a common macroeconomic driver along the South--North gradient
(corroborating, not established, evidence); and the recent deceleration,
consistent with the exhaustion of land still exposed to conversion in the
South, supported both by the elimination of competing hypotheses and by a
measured mechanism, the decline in the conversion rate as a region is
depleted. The text presents the methodology,
preliminary results, the limits of what the work claims, and the schedule
until the defense.

\textbf{Keywords}: land use and land cover; Cerrado; agricultural frontier;
Goiás; MapBiomas.
\end{otherlanguage*}
\end{resumo}
